\section{Consequences for Java and Other Languages}
\label{sec:impact}
The consequences of our work for Java are clear: there is an unsoundness that has existed since Java 5, and there is a way to avoid the unsoundness. But our results have bearing on similar languages as well, and comparing to languages similar to Java where no analogous situation arises helps articulate the places the idea is usefu.

\subsection{JVM Languages}


\paragraph{Flix}
Flix~\cite{madsen2016from} is a functional language for the JVM.
The Flix compiler monomorphizes polymorphic functions~\cite{madsen2023programming} (as Rust, C++, and MLton do) and has no subtyping.
The Flix construct closest to a local class is the \Q@new@ expression, which creates an object of an unnamed class that extends a Java class or implements a Java interface~\cite[\S14.4]{flix_book}
similar to Java's anonymous inner classes. Because a Flix program can not name such a class, casts or type-tests cannot be written against such types.

\paragraph{Kotlin}
Kotlin~\cite{kotlin_spec} already types a local class as a generic class.
This is the method \Q@m@ of Example~1 as a Kotlin function:

\begin{lstlisting}
open class Foo<X> { open fun x(): X = x() }
fun <X> m(other: Any, v: X): Foo<X> {
  class MyFoo : Foo<X>() { override fun x(): X = v }
  if (other !is MyFoo) { return MyFoo() } // error
  return other
}
\end{lstlisting}

Kotlin types the local class \Q@MyFoo@ as \Q@MyFoo<X>@.
Kotlin~2.4.21 rejects the type test \Q@other !is MyFoo@ with the error \texttt{cannot check for instance of erased type 'MyFoo<X (of fun <X> m)>'} (note that \texttt{MyFoo} has picked up an explicit type parameter).
If the cast \Q@other as MyFoo@ replaces the type test, Kotlin~2.4.21 accepts the function \Q@m@ with the warning \texttt{unchecked cast of 'Any' to 'MyFoo<X (of fun <X> m)>'}.
Kotlin~1.9.25 reports the same error and the same warning, in different words.
These are the error and the warning that funnelling brings to Example~1 (Section~\ref{sec:funnelling}).
Thus if we viewing Kotlin's implementation as an elaboration of the source language to an internal language using funnelling, our results can be interepreted as a proof that Kotlin's approach does fix this problem with method-local classes.

Kotlin conservatively adds to the type of a local class \emph{every} type parameter in scope: the type parameters of the enclosing function and the type parameters of the enclosing class.
It does this even for type parameters that the local class does not mention or otherwise depend on:
for a local class \Q@Plain@ with an empty body, declared in a function with a type parameter \Q@X@, Kotlin rejects the type test \Q@other is Plain@ with the same error.
This is granularity (2) of Section~\ref{sec:funnelling}, applied by the type system.
With funnelling, the type of a local class requires only the type variables that the local class funnels.
Thus \Q@Plain@ would need no funnel, and a type test against \Q@Plain@ is well typed.

\paragraph{Scala}
Perhaps the most prominent typed alternative JVM language is Scala.
This is Example~1 in Scala:

\begin{lstlisting}[language=Scala,numbers=left,xleftmargin=2em]
abstract class Foo[X] { def x: X }
object Broken {
  def m[X](other: Any, v: X): Foo[X] = {
    class MyFoo extends Foo[X] { def x: X = v }
    other match {
      case f: MyFoo => f (*\label{line:scalatest}*)
      case _ => new MyFoo
    }
  }
  def main(a: Array[String]): Unit = {
    val r1: Foo[Integer] = m[Integer]("anything", 5)
    val r2: Foo[String] = m[String](r1, "Hi")
    val s: String = r2.x (*\label{line:scalause}*)
    println(s.length)
  }
}
\end{lstlisting}

Scala~2.13.18 accepts this program with no warning, as Java accepts Example~1.
Scala~3.9.0 and Scala~3.10.1 accept this program with the warning \texttt{the type test for MyFoo cannot be checked at runtime because it's a local class} at Line~\ref{line:scalatest}.
Scala~3 reports this warning for a type test against any local class, also when the enclosing method has no generic parameters.
On the JVM, this program throws a \Q@ClassCastException@ at Line~\ref{line:scalause}: \Q@r2.x@ returns the \Q@Integer@ \Q@5@ where a \Q@String@ is expected.

There are also direct ahead-of-time compilers for both Scala and JVM bytecode.
Scala Native~\cite{shabalin2020} compiles Scala to machine code, and GraalVM Native Image~\cite{wimmer2019initialize} compiles JVM bytecode to machine code.
This program compiled by Scala Native~0.5.12, and Example~1 compiled by the Native Image of GraalVM for JDK~25, throw the same \Q@ClassCastException@.

Scala.js~\cite{scalajs_semantics} compiles Scala to JavaScript.
In Scala.js, a failing cast is undefined behaviour: a development build checks the cast, and a release build does not.
With Scala.js~1.22.0, a development build of this program throws an \Q@UndefinedBehaviorError@ at Line~\ref{line:scalause}.
In a release build, \Q@s@ is the JavaScript number \Q@5@, \Q@s.length@ reads the property \Q@length@ of that number, and the program prints \Q@undefined@.
Thus with Scala~2.13.18 and a release build of Scala.js, this program has no warning at compile time and no error at run time.
In Examples~1 and~2, the failure is detected at run time by a check that not every compilation target performs.

The key issues at work in the interaction of method-local inner classes and generics arise from the interaction of (1) nominal dynamic downcasts, (2) design choices about what runtime information about generic parameters is both persisted \emph{and checked} by casts, and (3) (not) tracking all dependencies on type parameters in the type system.
Essentially, Java and Scala both treat \Q@MyFoo@ as a sort of existential type --- that there exists some \Q@X@  such that
\Q@MyFoo@ extends \Q@Foo<X>@. However, treating \Q@MyFoo@ as a closed rather than generic type is equivalent
to assuming that there is only one existential for all instances --- associated with the declaration site --- rather than potentially
different existentials for different instances. This combined with the fact that the JVM is unable to dynamically check for different instantiations
leads to the unsoundness.
Our static solution effectively prohibits the de facto existential capture of the enclosing type parameter.


The decision to treat types like the \Q@MyFoo@ variants in our examples as closed existential types that depend on other choices of types is not
without parallels in other languages. Haskell and Rust, for example, include associated types~\cite{chakravarty2005associated}, where
a typeclass has a type as a member.
A yet closer analogue is Scala's path-dependent types~\cite{amin2016dot}, where types may be \emph{members} of an individual object, rather than associated with an individual class. The code below adapts the troublesome example method using Scala's type members:
\begin{lstlisting}[language=Scala]
trait Foo { type X; def x() : X; }
def m[Y](other: Any, newx: Y) : Foo2 = 
    class MyFoo extends Foo { type X = Y; override def x() : X = newx; }
    newx match { 
        case f: MyFoo => f
        case _ => MyFoo()
    }
\end{lstlisting}
This code compiles with the same warning that the cast cannot be checked because \texttt{MyFoo} is a local class, but in this case the warning is totally spurious:
compiling the corresponding analogue of Example~1's client code
\begin{lstlisting}[language=Scala]
def main(a: Array[String]) : Unit = {
    val r1: Foo = m[Integer]("anything", 5)
    val r2: Foo = m[String](r1, "Hi")
    val s: String = r2.x()
    println(s.length)
}
\end{lstlisting}
yields a static type error because the fact that \texttt{r2} (should) wrap a string is abstracted.
\todo[inline]{do we want to talk about this? kinda weird; would still be rejected by a naive mix of funnelling and DOTs. There's arguably a more general notion of ``capture dependency on enclosing generics'' that this particular example does safely}

\subsection{Other Languages}
Some other languages have features or feature sets very close to Java's method local classes, but avoid the problem with alternative language designs.

\paragraph{C\#}
The most similar language to Java at the source level makes very different design choices.
C\# has reified generics~\cite{kennedy2001design}: if \Q@o@ is a \Q@List<int>@, the type test \Q@o is List<string>@ is false and the cast \Q@(List<string>)o@ throws an \Q@InvalidCastException@.
C\# has no local classes.
The closest C\# constructs are anonymous types and lambdas, which the C\# compiler translates to generated classes.
In a generic method \Q@M<X>@ with a parameter \Q@x@ of type \Q@X@, the .NET SDK~10.0.404 translates \Q@new { v= x }@ and \Q@()=>x@ to instances of generated classes that are generic in \Q@X@.
Thus the objects created by \Q@M<int>@ and by \Q@M<string>@ have different classes: this is granularity (2) of Section~\ref{sec:funnelling}, in the type system and at run time.
A C\# program can not name a class that was generated by the compiler,\footnote{The name itself is only chosen during compilation, and the C\# compiler exploits the fact that the .NET CLR admits type names that are not valid C\# source level identifiers, such as \lstinline|<>f__AnonymousType0|.} so a C\# program can not write a type test against a generated class directly.
As a result of such decisions, most languages that are potentially vulnerable to this issue are JVM languages.

\paragraph{OCaml}
Since version 4.04, OCaml has included local exceptions, where an exception may be declared local to a function.
These exceptions exhibit category 3 in the taxomony of \S\ref{sec:funnelling}, where each dynamic invocation yields a distinct identity. The documentation~\cite[\S7]{ocamldocs} gives the following example to demonstrate this:
\begin{lstlisting}[language=Caml]
# let gen () = let exception A in A

  let () = assert(gen () = gen ());;
Exception: Assert_failure ("expr.etex", 3, 9).
\end{lstlisting}
OCaml already tracks enough information for a cast between such types to fail, but it does not include casts, type-cases, or similar features (even the object-oriented part of OCaml lacks casts).
